\documentclass[uplatex,dvipdfmx,report]{jlreq}
\usepackage[fleqn,tbtags]{mathtools}
\usepackage{jlreq-deluxe}
\usepackage{url}
\usepackage[uplatex,deluxe]{otf}
\usepackage[noalphabet]{pxchfon}
\usepackage{stix2}
\usepackage{amsmath}
\usepackage{listings}
\usepackage{xcolor}
\usepackage{geometry}
\usepackage{graphicx}
\usepackage{float}
\usepackage{longtable}
\usepackage{multirow}

\geometry{margin=25mm}

\lstset{
  basicstyle=\ttfamily\small,
  keywordstyle=\color{blue},
  commentstyle=\color{gray},
  stringstyle=\color{red},
  frame=single,
  breaklines=true,
  numbers=left,
  numberstyle=\tiny,
  language=C++
}

\title{Unvanilla Industry}
\author{Eval02}

\begin{document}

\maketitle

\chapter{フォルダ構造}
\begin{table}[H]
    \centering
    \caption{scripts}
    \begin{tabular}{cc}
        \hline
        ファイル名 & 機能 \\
        \hline
        \hline
        units & ユニットのデータを保存する \\
        main.js & メインの処理を行う \\
        utils.js & 全体で利用するfunctionを所有 \\
        slotClass.js & スロットに関するclassを所有 \\
        powerNetworkClass.js & 電気網に関するclassを所有 \\
        unitClass.js & ユニット系のclassの継承関係を管理する \\
        \hline
    \end{tabular}
\end{table}

\begin{table}[H]
    \centering
    \caption{units}
    \begin{tabular}{c}
        \hline
        ユニット\\
        \hline
        \hline
        powerCable \\
        minerMk1 \\
        coalGenerator \\
        furnaceMachine \\
        conveyorMk1 \\
        pipeMk1 \\
        \hline
    \end{tabular}
\end{table}

\begin{table}[H]
    \centering
    \caption{ユニット内部}
    \begin{tabular}{cc}
        \hline
        ファイル名 & 機能 \\
        \hline
        \hline
        class.js & ユニットの機能を処理するclassを所有 \\
        data.js & ユニットの初期情報を設定するデータを所有 \\
        form.js & ユニットのUIを所有 \\
        recipe.js & ユニットが扱うレシピを所有 \\
        structure.js & ユニットのブロック構造を所有 \\
        \hline
    \end{tabular}
\end{table}

\chapter{scripts直下ファイル内部}

\section{main.js}

\begin{table}[H]
    \centering
    \caption{所有オブジェクト}
    \begin{tabular}{cc}
        \hline
        オブジェクト名 & 機能 \\
        \hline
        \hline
        unitUuidMap & ユニットのuuidとユニットのインスタンスを結び付けるmap \\
        slotUuidMap & スロットのuuidとスロットのインスタンスを結び付けるmap \\
        slotPosMap & スロットの座標とスロットのインスタンスを結び付けるmap \\
        powerNetworkMap & パワーネットワークのuuidとインスタンスを結び付けるmap \\
        \hline
    \end{tabular}
\end{table}

\section{slotClass.js}

\subsection{class一覧}
\begin{table}[H]
    \centering
    \caption{所有class}
    \begin{tabular}{ccc}
        \hline
        class名 & 継承 & 機能 \\
        \hline
        \hline
        Slot & - & スロット系classの共通機能を所有する \\
        IOSlot & Slot & アイテム、液体を扱うスロットの機能を所有する \\
        ElectrodeSlot & Slot & 電極を扱うスロットの機能を所有する \\
        \hline
    \end{tabular}
\end{table}

\subsection{Slot}

Slot系Classの基本的な役割を担う．

\begin{table}[H]
    \centering
    \caption{必要な機能}
    \begin{tabular}{cc}
        \hline
        機能 & メソッド名 \\
        \hline
        \hline
        接続先を探索する & searchConnect \\
        \hline
    \end{tabular}
\end{table}

\subsubsection{constructor}

\begin{table}[H]
    \centering
    \caption{入力}
    \begin{tabular}{cc}
        \hline
        変数名 & 機能 \\
        \hline
        \hline
        parentId & 親ユニットのuuid \\
        index & 親ユニットからの自身のスロット番号 \\
        \hline
    \end{tabular}
\end{table}

\begin{table}[H]
    \centering
    \caption{変数}
    \begin{tabular}{cc}
        \hline
        所有変数 & 機能 \\
        \hline
        \hline
        parentId & 親ユニットのuuid \\
        parentUnit & 親ユニット \\
        index & 自身のスロット番号 \\
        slotData & スロットの初期セットデータ \\
        pos & スロットの絶対座標 \\
        face & スロットの面 \\
        dimension & スロットのディメンション \\
        connectId & 接続先スロットのuuid \\
        connectSlot & 接続先スロット \\
        \hline
    \end{tabular}
\end{table}

\subsubsection{searchConnect}

\begin{table}[H]
    \centering
    \caption{入力}
    \begin{tabular}{cc}
        \hline
        変数名 & 機能 \\
        \hline
        \hline
        なし & - \\
        \hline
    \end{tabular}
\end{table}

pos，faceより接続先座標を計算．\\
接続先座標をもとにMain.slotPosMapよりslotを探索．\\
connectIdに接続先slotのuuidを代入し，接続先slotのconnectIdに自分のuuidを代入する．

\subsection{IOSlot}

Slotを継承．アイテム及び液体の入出力スロットの役割を担う．

\begin{table}[H]
    \centering
    \caption{必要な機能}
    \begin{tabular}{cc}
        \hline
        機能 & メソッド名 \\
        \hline
        \hline
        接続先を探索する & searchConnect \\
        inputスロットにアイテムを入力する & inputItem \\
        outputスロットからアイテムを出力する & outputItem \\
        \hline
    \end{tabular}
\end{table}

\subsubsection{constructor}

\begin{table}[H]
    \centering
    \caption{入力}
    \begin{tabular}{cc}
        \hline
        変数名 & 機能 \\
        \hline
        \hline
        Slotの継承 & - \\
        \hline
    \end{tabular}
\end{table}

\begin{table}[H]
    \centering
    \caption{変数}
    \begin{tabular}{cc}
        \hline
        所有変数 & 機能 \\
        \hline
        \hline
        Slotの継承 & - \\
        content & スロットが所有するコンテンツ \\
        maxAmount & スロットが所有できるコンテンツの最大量 \\
        \hline
    \end{tabular}
\end{table}

\subsubsection{inputItem}

\begin{table}[H]
    \centering
    \caption{入力}
    \begin{tabular}{cc}
        \hline
        変数名 & 機能 \\
        \hline
        \hline
        なし & - \\
        \hline
    \end{tabular}
\end{table}

接続先スロットからアイテムを移動させる．

\subsubsection{outputItem}

\begin{table}[H]
    \centering
    \caption{入力}
    \begin{tabular}{cc}
        \hline
        変数名 & 機能 \\
        \hline
        \hline
        なし & - \\
        \hline
    \end{tabular}
\end{table}

接続先スロットにアイテムを移動する．

\subsection{ElectrodeSlot}

Slotを継承．電極スロットの役割を担う．

\begin{table}[H]
    \centering
    \caption{必要な機能}
    \begin{tabular}{cc}
        \hline
        機能 & メソッド名 \\
        \hline
        \hline
        接続先を探索する & searchConnect \\
        \hline
    \end{tabular}
\end{table}

\subsubsection{constructor}

\begin{table}[H]
    \centering
    \caption{入力}
    \begin{tabular}{cc}
        \hline
        変数名 & 機能 \\
        \hline
        \hline
        Slotの継承 & - \\
        \hline
    \end{tabular}
\end{table}

\begin{table}[H]
    \centering
    \caption{変数}
    \begin{tabular}{cc}
        \hline
        所有変数 & 機能 \\
        \hline
        \hline
        Slotの継承 & - \\
        electrodeType & 電気の発電，消費，接続を判別（generate/consume/connect） \\
        powerPerMinute & 分間電力量（generate: 発電量，consume: 消費量，connect: 0） \\
        \hline
    \end{tabular}
\end{table}

\section{powerNetworkClass.js}

\subsection{class一覧}

\begin{table}[H]
    \centering
    \caption{所有class}
    \begin{tabular}{ccc}
        \hline
        class名 & 継承 & 機能 \\
        \hline
        \hline
        PowerNetwork & - & 電気網の管理を行う \\
        \hline
    \end{tabular}
\end{table}

\subsection{PowerNetwork}

\begin{table}[H]
    \centering
    \caption{必要な機能}
    \begin{tabular}{cc}
        \hline
        機能 & メソッド名 \\
        \hline
        \hline
        電力の更新 & updatePower \\
        メンバーの追加 & addMember \\
        メンバーの削除 & removeMember \\
        \hline
    \end{tabular}
\end{table}

\subsubsection{constructor}

\begin{table}[H]
    \centering
    \caption{変数}
    \begin{tabular}{cc}
        \hline
        所有変数 & 機能 \\
        \hline
        \hline
        uuid & ネットワークのuuid \\
        memberIds & 所属ユニットのuuidリスト \\
        isPowered & 電力供給状態 \\
        generationPerMinute & 分間発電量 \\
        consumptionPerMinute & 分間消費量 \\
        \hline
    \end{tabular}
\end{table}

\subsubsection{updatePower}

全メンバーの発電量・消費量を集計し，isPoweredを更新する．\\
発電量が消費量を上回る場合はisPoweredをtrueに，下回る場合はfalseにする．

\subsubsection{addMember}

\begin{table}[H]
    \centering
    \caption{入力}
    \begin{tabular}{cc}
        \hline
        変数名 & 機能 \\
        \hline
        \hline
        unitId & 追加するユニットのuuid \\
        \hline
    \end{tabular}
\end{table}

memberIdsにunitIdを追加しupdatePowerを呼ぶ．

\subsubsection{removeMember}

\begin{table}[H]
    \centering
    \caption{入力}
    \begin{tabular}{cc}
        \hline
        変数名 & 機能 \\
        \hline
        \hline
        unitId & 削除するユニットのuuid \\
        \hline
    \end{tabular}
\end{table}

memberIdsからunitIdを削除しupdatePowerを呼ぶ．\\
memberIdsが空になった場合はpowerNetworkMapから自身を削除する．

\section{unitClass.js}

\subsection{class一覧}
\begin{table}[H]
    \centering
    \caption{所有class}
    \begin{tabular}{ccc}
        \hline
        class名 & 継承 & 機能 \\
        \hline
        \hline
        Unit & - & 全ユニットの共通機能を所有する \\
        Machine & Unit & 機械系ユニットの共通機能を所有する \\
        Extraction & Machine & 資源採取系ユニットの共通機能を所有する \\
        Generator & Machine & 発電機系ユニットの共通機能を所有する \\
        Production & Machine & 生産系ユニットの共通機能を所有する \\
        Transport & Unit & 輸送系ユニットの共通機能を所有する \\
        Conveyor & Transport & コンベヤ系ユニットの共通機能を所有する \\
        Pipe & Transport & パイプ系ユニットの共通機能を所有する \\
        Connector & Unit & コネクタ系ユニットの共通機能を所有する \\
        Cable & Connector & ケーブル系ユニットの共通機能を所有する \\
        \hline
    \end{tabular}
\end{table}

\subsection{Unit}

全ユニットの基底クラス．

\begin{table}[H]
    \centering
    \caption{必要な機能}
    \begin{tabular}{cc}
        \hline
        機能 & メソッド名 \\
        \hline
        \hline
        設置時のブロック配置 & setStructure \\
        接続先の探索 & searchConnect \\
        アニメーション & animation \\
        破壊時の後処理 & onDestroy \\
        シリアライズ & serialize \\
        デシリアライズ & fromJSON \\
        \hline
    \end{tabular}
\end{table}

\subsubsection{constructor}

\begin{table}[H]
    \centering
    \caption{入力}
    \begin{tabular}{cc}
        \hline
        入力 & 機能 \\
        \hline
        \hline
        pos & ユニットの絶対座標 \\
        dimension & ユニットのディメンション \\
        direction & ユニットの設置方向 \\
        \hline
    \end{tabular}
\end{table}

\begin{table}[H]
    \centering
    \caption{変数}
    \begin{tabular}{cc}
        \hline
        変数名 & 機能 \\
        \hline
        \hline
        uuid & ユニット別のuuid \\
        pos & ユニットの絶対座標 \\
        dimension & ユニットのディメンション \\
        direction & ユニットの設置方向 \\
        slots & スロットのリスト \\
        \hline
    \end{tabular}
\end{table}

\subsection{Machine}

Unitを継承．機械系ユニットの共通機能を担う．

\begin{table}[H]
    \centering
    \caption{必要な機能}
    \begin{tabular}{cc}
        \hline
        機能 & メソッド名 \\
        \hline
        \hline
        電力ネットワークへの参加 & joinPowerNetwork \\
        入力スロットの変化検知 & inputChange \\
        出力スロットの変化検知 & outputChange \\
        レシピ選択UIの実行 & selectRecipe \\
        レシピの処理 & doRecipe \\
        tickの処理 & tick \\
        \hline
    \end{tabular}
\end{table}

\subsubsection{constructor}

\begin{table}[H]
    \centering
    \caption{入力}
    \begin{tabular}{cc}
        \hline
        入力 & 機能 \\
        \hline
        \hline
        Unitの継承 & - \\
        \hline
    \end{tabular}
\end{table}

\begin{table}[H]
    \centering
    \caption{変数}
    \begin{tabular}{cc}
        \hline
        変数名 & 機能 \\
        \hline
        \hline
        Unitの継承 & - \\
        status & 機械の動作状態（IDLE/RUNNING/WAITING\_OUTPUT/BLACKOUT） \\
        powerNetworkId & 所属する電力ネットワークのuuid \\
        currentRecipeId & 現在選択中のレシピのid \\
        cycleStartTick & 現在のサイクルの開始tick \\
        cycleEndTick & 現在のサイクルの終了tick \\
        \hline
    \end{tabular}
\end{table}

\subsection{Extraction}

Machineを継承．資源採取系ユニットの共通機能を担う．

\begin{table}[H]
    \centering
    \caption{必要な機能}
    \begin{tabular}{cc}
        \hline
        機能 & メソッド名 \\
        \hline
        \hline
        blockSensorの更新 & updateBlockSensor \\
        \hline
    \end{tabular}
\end{table}

\subsubsection{constructor}

\begin{table}[H]
    \centering
    \caption{入力}
    \begin{tabular}{cc}
        \hline
        入力 & 機能 \\
        \hline
        \hline
        Machineの継承 & - \\
        \hline
    \end{tabular}
\end{table}

\begin{table}[H]
    \centering
    \caption{変数}
    \begin{tabular}{cc}
        \hline
        変数名 & 機能 \\
        \hline
        \hline
        Machineの継承 & - \\
        detectedBlockId & blockSensorが検知したブロックのid \\
        \hline
    \end{tabular}
\end{table}

\subsection{Generator}

Machineを継承．発電機系ユニットの共通機能を担う．

\begin{table}[H]
    \centering
    \caption{必要な機能}
    \begin{tabular}{cc}
        \hline
        機能 & メソッド名 \\
        \hline
        \hline
        ブレーカーの操作 & toggleBreaker \\
        \hline
    \end{tabular}
\end{table}

\subsubsection{constructor}

\begin{table}[H]
    \centering
    \caption{入力}
    \begin{tabular}{cc}
        \hline
        入力 & 機能 \\
        \hline
        \hline
        Machineの継承 & - \\
        \hline
    \end{tabular}
\end{table}

\begin{table}[H]
    \centering
    \caption{変数}
    \begin{tabular}{cc}
        \hline
        変数名 & 機能 \\
        \hline
        \hline
        Machineの継承 & - \\
        breakerActive & ブレーカーの状態 \\
        \hline
    \end{tabular}
\end{table}

\subsection{Production}

Machineを継承．生産系ユニットの共通機能を担う．

\subsubsection{constructor}

\begin{table}[H]
    \centering
    \caption{入力}
    \begin{tabular}{cc}
        \hline
        入力 & 機能 \\
        \hline
        \hline
        Machineの継承 & - \\
        \hline
    \end{tabular}
\end{table}

\begin{table}[H]
    \centering
    \caption{変数}
    \begin{tabular}{cc}
        \hline
        変数名 & 機能 \\
        \hline
        \hline
        Machineの継承 & - \\
        \hline
    \end{tabular}
\end{table}

\subsection{Transport}

Unitを継承．輸送系ユニットの共通機能を担う．

\begin{table}[H]
    \centering
    \caption{必要な機能}
    \begin{tabular}{cc}
        \hline
        機能 & メソッド名 \\
        \hline
        \hline
        アイテムの輸送 & transportItem \\
        tickの処理 & tick \\
        \hline
    \end{tabular}
\end{table}

\subsubsection{constructor}

\begin{table}[H]
    \centering
    \caption{入力}
    \begin{tabular}{cc}
        \hline
        入力 & 機能 \\
        \hline
        \hline
        Unitの継承 & - \\
        \hline
    \end{tabular}
\end{table}

\begin{table}[H]
    \centering
    \caption{変数}
    \begin{tabular}{cc}
        \hline
        変数名 & 機能 \\
        \hline
        \hline
        Unitの継承 & - \\
        currentContent & 現在保持しているコンテンツ \\
        nextContent & 次tickに反映するコンテンツ（ダブルバッファリング用） \\
        \hline
    \end{tabular}
\end{table}

\subsection{Conveyor}

Transportを継承．コンベヤ系ユニットの共通機能を担う．

\subsubsection{constructor}

\begin{table}[H]
    \centering
    \caption{入力}
    \begin{tabular}{cc}
        \hline
        入力 & 機能 \\
        \hline
        \hline
        Transportの継承 & - \\
        \hline
    \end{tabular}
\end{table}

\begin{table}[H]
    \centering
    \caption{変数}
    \begin{tabular}{cc}
        \hline
        変数名 & 機能 \\
        \hline
        \hline
        Transportの継承 & - \\
        \hline
    \end{tabular}
\end{table}

\subsection{Pipe}

Transportを継承．パイプ系ユニットの共通機能を担う．

\subsubsection{constructor}

\begin{table}[H]
    \centering
    \caption{入力}
    \begin{tabular}{cc}
        \hline
        入力 & 機能 \\
        \hline
        \hline
        Transportの継承 & - \\
        \hline
    \end{tabular}
\end{table}

\begin{table}[H]
    \centering
    \caption{変数}
    \begin{tabular}{cc}
        \hline
        変数名 & 機能 \\
        \hline
        \hline
        Transportの継承 & - \\
        pumpId & 影響を受けるポンプのuuid \\
        role & ポンプに対する役割（source/dest） \\
        \hline
    \end{tabular}
\end{table}

\subsection{Connector}

Unitを継承．コネクタ系ユニットの共通機能を担う．

\subsubsection{constructor}

\begin{table}[H]
    \centering
    \caption{入力}
    \begin{tabular}{cc}
        \hline
        入力 & 機能 \\
        \hline
        \hline
        Unitの継承 & - \\
        \hline
    \end{tabular}
\end{table}

\begin{table}[H]
    \centering
    \caption{変数}
    \begin{tabular}{cc}
        \hline
        変数名 & 機能 \\
        \hline
        \hline
        Unitの継承 & - \\
        \hline
    \end{tabular}
\end{table}

\subsection{Cable}

Connectorを継承．電力ケーブルの共通機能を担う．

\subsubsection{constructor}

\begin{table}[H]
    \centering
    \caption{入力}
    \begin{tabular}{cc}
        \hline
        入力 & 機能 \\
        \hline
        \hline
        Connectorの継承 & - \\
        \hline
    \end{tabular}
\end{table}

\begin{table}[H]
    \centering
    \caption{変数}
    \begin{tabular}{cc}
        \hline
        変数名 & 機能 \\
        \hline
        \hline
        Connectorの継承 & - \\
        \hline
    \end{tabular}
\end{table}

\chapter{各ユニット内部}

\section{基本的なデータ構造}

\begin{table}[H]
    \centering
    \caption{データ構造}
    \begin{tabular}{cc}
        \hline
        プロパティ名 & 機能 \\
        \hline
        \hline
        typeId & ユニットの種別ごとのid \\
        category & ユニットのカテゴリ \\
        subCategory & ユニットのサブカテゴリ \\
        inputSlots & 入力スロットの設定 \\
        outputSlots & 出力スロットの設定 \\
        electrodeSlots & 電極スロットの設定 \\
        \hline
    \end{tabular}
\end{table}

\begin{table}[H]
    \centering
    \caption{入出力スロット}
    \begin{tabular}{cc}
        \hline
        プロパティ名 & 機能 \\
        \hline
        \hline
        slotType & スロットの種別ごとのid \\
        ioType & 入力，出力を判別 \\
        localPos & スロットの相対座標 \\
        face & スロットの面 \\
        contentType & 扱うコンテンツの種別 \\
        maxAmount & スロットの最大スタック数 \\
        \hline
    \end{tabular}
\end{table}

\begin{table}[H]
    \centering
    \caption{電極スロット}
    \begin{tabular}{cc}
        \hline
        プロパティ名 & 機能 \\
        \hline
        \hline
        slotType & スロットの種別ごとのid \\
        electrodeType & 電気の発電，消費，接続を判別 \\
        localPos & スロットの相対座標 \\
        face & スロットの面 \\
        \hline
    \end{tabular}
\end{table}

\section{minerMk1}

\subsection{minerMk1Data.js}

\begin{table}[H]
    \centering
    \caption{データ構造}
    \begin{tabular}{cc}
        \hline
        プロパティ名 & 値 \\
        \hline
        \hline
        typeId & minerMk1 \\
        category & MACHINE \\
        subCategory & EXTRACTION \\
        inputSlots & 0 \\
        outputSlots & 1 \\
        electrodeSlots & 1 \\
        \hline
    \end{tabular}
\end{table}

\subsection{minerMk1Class.js}

Extractionを継承．

\begin{table}[H]
    \centering
    \caption{必要な機能}
    \begin{tabular}{cc}
        \hline
        機能 & メソッド名 \\
        \hline
        \hline
        設置時のブロック配置 & setStructure \\
        入力スロットの変化検知 & inputChange \\
        出力スロットの変化検知 & outputChange \\
        レシピ選択UIの実行 & selectRecipe \\
        レシピの処理 & doRecipe \\
        アニメーション & animation \\
        \hline
    \end{tabular}
\end{table}

\subsection{minerMk1Recipe.js}

\begin{table}[H]
    \centering
    \caption{レシピ一覧}
    \begin{tabular}{ccccc}
        \hline
        レシピid & 入力 & 出力 & 消費電力(w/m) & 処理時間(tick) \\
        \hline
        \hline
        minerIron & 鉄鉱石ブロック検知 & 鉄鉱石x1 & 30 & 20 \\
        minerCoal & 石炭鉱石ブロック検知 & 石炭鉱石x1 & 30 & 20 \\
        \hline
    \end{tabular}
\end{table}

\section{coalGenerator}

\subsection{coalGeneratorData.js}

\begin{table}[H]
    \centering
    \caption{データ構造}
    \begin{tabular}{cc}
        \hline
        プロパティ名 & 値 \\
        \hline
        \hline
        typeId & coalGenerator \\
        category & MACHINE \\
        subCategory & GENERATOR \\
        inputSlots & 1 \\
        outputSlots & 0 \\
        electrodeSlots & 1 \\
        \hline
    \end{tabular}
\end{table}

\subsection{coalGeneratorClass.js}

Generatorを継承．

\begin{table}[H]
    \centering
    \caption{必要な機能}
    \begin{tabular}{cc}
        \hline
        機能 & メソッド名 \\
        \hline
        \hline
        設置時のブロック配置 & setStructure \\
        入力スロットの変化検知 & inputChange \\
        出力スロットの変化検知 & outputChange \\
        レシピ選択UIの実行 & selectRecipe \\
        レシピの処理 & doRecipe \\
        アニメーション & animation \\
        ブレーカーの操作 & toggleBreaker \\
        \hline
    \end{tabular}
\end{table}

\subsection{coalGeneratorRecipe.js}

\begin{table}[H]
    \centering
    \caption{レシピ一覧}
    \begin{tabular}{ccccc}
        \hline
        レシピid & 入力 & 出力 & 発電量(w/m) & 処理時間(tick) \\
        \hline
        \hline
        genCoal & 石炭x1 & なし & 300 & 20 \\
        \hline
    \end{tabular}
\end{table}

\section{furnaceMachine}

\subsection{furnaceMachineData.js}

\begin{table}[H]
    \centering
    \caption{データ構造}
    \begin{tabular}{cc}
        \hline
        プロパティ名 & 値 \\
        \hline
        \hline
        typeId & furnaceMachine \\
        category & MACHINE \\
        subCategory & PRODUCTION \\
        inputSlots & 1 \\
        outputSlots & 1 \\
        electrodeSlots & 1 \\
        \hline
    \end{tabular}
\end{table}

\subsection{furnaceMachineClass.js}

Productionを継承．

\begin{table}[H]
    \centering
    \caption{必要な機能}
    \begin{tabular}{cc}
        \hline
        機能 & メソッド名 \\
        \hline
        \hline
        設置時のブロック配置 & setStructure \\
        入力スロットの変化検知 & inputChange \\
        出力スロットの変化検知 & outputChange \\
        レシピ選択UIの実行 & selectRecipe \\
        レシピの処理 & doRecipe \\
        アニメーション & animation \\
        \hline
    \end{tabular}
\end{table}

\subsection{furnaceMachineRecipe.js}

\begin{table}[H]
    \centering
    \caption{レシピ一覧}
    \begin{tabular}{ccccc}
        \hline
        レシピid & 入力 & 出力 & 消費電力(w/m) & 処理時間(tick) \\
        \hline
        \hline
        smeltIron & 鉄鉱石x1 & 鉄インゴットx1 & 30 & 20 \\
        smeltCoal & 石炭鉱石x1 & 石炭x1 & 30 & 20 \\
        \hline
    \end{tabular}
\end{table}

\section{conveyorMk1}

\subsection{conveyorMk1Data.js}

\begin{table}[H]
    \centering
    \caption{データ構造}
    \begin{tabular}{cc}
        \hline
        プロパティ名 & 値 \\
        \hline
        \hline
        typeId & conveyorMk1 \\
        category & TRANSPORT \\
        subCategory & CONVEYOR \\
        inputSlots & 1 \\
        outputSlots & 1 \\
        electrodeSlots & 0 \\
        speedPerMinute & 60 \\
        \hline
    \end{tabular}
\end{table}

\subsection{conveyorMk1Class.js}

Conveyorを継承．

\begin{table}[H]
    \centering
    \caption{必要な機能}
    \begin{tabular}{cc}
        \hline
        機能 & メソッド名 \\
        \hline
        \hline
        設置時のブロック配置 & setStructure \\
        アイテムの輸送 & transportItem \\
        次tickへの反映 & applyNext \\
        アニメーション & animation \\
        \hline
    \end{tabular}
\end{table}

\end{document}